\section*{Coda: spectra inherit every earlier choice}
\addcontentsline{toc}{section}{Coda: spectra inherit every earlier choice}

A spectrum inherits the meaning of the graph from which it was built.  Gauge projection,
scalarization, symmetrization, stationary weighting, normalization, context averaging, and null
subtraction all change the operator being diagonalized.  Reporting those choices is therefore
part of the spectral result.

Used with that discipline, spectral analysis becomes a genuine multiscale extension of the
operator framework.  Scalar modes describe positional smoothness, signed modes encode frustrated
polarity, connection modes probe transported categorical coherence, and directed decompositions
separate reversible flow from circulation.  Comparing their projectors can reveal organization
that no single edge statistic exposes.

The word ``mode'' carries a physical temptation that must be resisted until a dynamics has been
declared.  An eigenvector of a scalar graph Laplacian is a field that is smooth relative to chosen
edge weights.  It becomes a vibrational mode only when the operator is a stiffness matrix with a
mass metric; it becomes a relaxation mode only when the operator generates a justified dissipative
dynamics; and it becomes a diffusion mode only when the state variable and transition rates have
physical meaning.  Without those structures, low frequency means graph smoothness and nothing
more.

The distinction is sharper for the other spectra.  A connection-Laplacian mode transports amino-
acid contrast directions, not atomic displacement vectors.  A signed mode records coherence or
frustration under a declared scalar contraction, not automatically favorable or unfavorable
thermodynamic coupling.  Entropy production of an attention route quantifies irreversibility of a
computational Markov chain; it is literal thermodynamic entropy production only if that chain has
been identified with a physical stochastic process.  Spectral resemblance to sectors, allosteric
pathways, or collective motions is therefore a hypothesis to test against structural ensembles,
perturbation responses, or time-resolved dynamics, not a conclusion supplied by diagonalization.

The remaining question is constructive.  If the desired interaction object is now clear, can a
trainable architecture realize it while preserving the measurability conditions that gave it
meaning?  Sparse support, dynamic routing, endpoint exclusion, persistent edge state, and
efficient execution jointly define a genuine design problem.  The next part turns from
interpretation to realizability.
